\Needspace{14\baselineskip}
\section{Source promises, normalizations, and provenance}
\label{sec:source-normalizations}

The source derivations follow the order in which the paper uses them:
the bounded-degree independent-set promise, the multipartite UGC
promise, the global SSE promise, and the NP-hardness provenance and
parity normalization.

\subsection{Degree reduction for the UGC independent-set source}
\label{app:bounded-degree-ugc-source}

We derive Theorem~\ref{thm:bounded-degree-kr-source} from the
weighted density theorem of Austrin, Khot, and Safra. Their later
bounded-degree reduction uses random sampling
\cite[Section~4, Steps~2--3]{AKS2011}. We instead obtain bounded
degree deterministically by spectral sparsification. We state the
cited density theorem first, then derive the loop removal, vertex
unweighting, and degree bounds needed here.

\subsubsection{The AKS density theorem}

A bipartite \emph{Unique Games instance}
$\Lambda=(U,V,E_\Lambda,\Pi,[\ell])$ consists of a bipartite
multigraph with left vertices $U$, right vertices $V$, and a label
set $[\ell]$. Each edge $e=uv$ has a permutation $\pi_e\in\Pi$
and requires the label at $v$ to be $\pi_e$ applied to the label at
$u$. Write $\OPT(\Lambda)$ for the maximum fraction
of satisfied constraints, counting parallel constraints separately.
Write $\operatorname{VOPT}(\Lambda)$ for the largest fraction
$|V_0|/|V|$ for which a labeling satisfies every constraint incident
to $V_0\subseteq V$. The instance is \emph{regular} if all left
vertices have the same degree and all right vertices have the same
degree. The two degrees may differ.

\begin{theorem}[Regular Unique Games source
{\cite[Conjecture~2.4]{AKS2011} and \cite{KhotRegev}}]
\label{thm:regular-ug-source}
For every fixed $\gamma>0$, there is a fixed alphabet size $\ell$
such that it is UGC-hard to distinguish regular bipartite Unique
Games instances $\Lambda$ with alphabet $[\ell]$ satisfying the
following cases.
\begin{description}
\item[YES:] $\operatorname{VOPT}(\Lambda)\geq1-\gamma$.
 There is a set $V_0\subseteq V$ of at least $(1-\gamma)|V|$ right
 vertices and one labeling satisfying every constraint incident
 to $V_0$.
\item[NO:] $\OPT(\Lambda)\leq\gamma$.
 Every labeling satisfies at most a $\gamma$ fraction of the
 constraints.
\end{description}
\end{theorem}

For a vertex- and edge-weighted graph $G$, let $a_v$ denote a
vertex weight and $w_e$ an edge weight, normalized so that each
of the two total weights is one. Put
$a(S)=\sum_{v\in S}a_v$ and let $e_G(S)$ be the total weight of
edges with both endpoints in $S$. We allow loops, counted once in
$e_G(S)$ and twice in the weighted degree $\deg_G(v)$.
For $\beta\in[0,1]$ and a real $\delta$, the graph is
\emph{$(\beta,\delta)$-dense} if $a(S)\geq\beta$ implies
$e_G(S)\geq\delta$.

Let $\Psi$ be the standard normal cumulative distribution function.
For $-1<\rho<1$, let $X,Y$ be jointly normal with mean zero,
variance one, and covariance $\rho$, and define
\[
 \Gamma_\rho(s)=
 \Pr[X\leq\Psi^{-1}(s),\ Y\leq\Psi^{-1}(s)]
 \qquad(0<s<1).
\]
Set $\Gamma_\rho(0)=0$ and $\Gamma_\rho(1)=1$.
Fix a bias $q\in(0,1/2)$ and an analytic error
$\varepsilon_{\mathrm{dens}}>0$, and put $\rho=-q/(1-q)$.

\Needspace{11\baselineskip}
\begin{theorem}[AKS density theorem {\cite[Theorem~3.1]{AKS2011}}]
\label{thm:aks-density-source}
For all sufficiently small $\gamma>0$, an algorithm maps a Unique
Games instance $\Lambda$ to a vertex- and edge-weighted graph $G$
with the following properties.
\begin{description}
\item[YES:] If $\operatorname{VOPT}(\Lambda)\geq1-\gamma$, then
 $G$ has an independent set of vertex weight at least $q-\gamma$.
\item[NO:] If $\OPT(\Lambda)\leq\gamma$ and
 $\Lambda$ is regular, then $G$ is
 $(\beta,\Gamma_\rho(\beta)-\varepsilon_{\mathrm{dens}})$-dense
 for every $\beta\in[0,1]$.
\end{description}
Vertex weights are proportional to weighted degrees.
The running time is polynomial in $|U|,|V|,|E_\Lambda|$ and
exponential in $\ell$.
\end{theorem}

For a regular input, the proportionality gives
$\deg_G(v)=2a_v$: the weighted degrees sum to twice the total edge
weight, whereas the vertex weights sum to one. Equivalently, each
endpoint marginal of the AKS edge distribution is $a_v$, with a
loop contributing twice. Thus the NO conclusion, applied with
$\beta=a(S)$, is
\[
 e_G(S)\geq\Gamma_\rho(a(S))-\varepsilon_{\mathrm{dens}}
 \qquad(S\subseteq V(G)).
\]

We will use positivity, monotonicity, and convexity of $\Gamma_\rho$.
For completeness, these follow directly from its definition. The
joint normal density is positive, so $\Gamma_\rho(s)>0$ for
$0<s<1$. Differentiating the bivariate distribution function along
its diagonal, using symmetry of $X,Y$ and the conditional normal
law, gives
\[
 \Gamma_\rho'(s)=
 2\Psi\!\left(\sqrt{\frac{1-\rho}{1+\rho}}\,\Psi^{-1}(s)\right).
\]
The right-hand side is nonnegative and increasing in $s$.
Consequently $\Gamma_\rho$ is nondecreasing and convex, including
at the endpoints by continuity.

\subsubsection{Removing loops and vertex weights}

Fix a rational bias $q$ from now on. The AKS sampling construction
can produce loops. We remove them before replacing vertices by
independent copies, so that the density argument can use
independent choices of distinct original vertices. Let
$\varepsilon_{\mathrm{loop}}>0$ be a fixed loop-loss budget.
The following alphabet padding makes the total loop weight smaller
than this budget.

\begin{lemma}[Removing loops by alphabet padding]\label{lem:aks-loop-removal}
For every fixed $\varepsilon_{\mathrm{loop}}>0$, a fixed alphabet padding followed by the
AKS construction and deletion of all loops gives a loopless graph $G^-$
with vertex weights $a_v$ of total one, total edge weight at most
one, and weighted degrees at most $2a_v$.
Its YES independent-set weight is at least $q-\gamma$, and in the
NO case
\begin{equation}\label{eq:aks-loop-free-density}
 e_{G^-}(S)\ge\Gamma_\rho(a(S))-\varepsilon_{\mathrm{dens}}-\varepsilon_{\mathrm{loop}}
 \qquad(S\subseteq V(G^-)).
\end{equation}
The reduction is deterministic and polynomial for fixed parameters.
\end{lemma}

\begin{proof}
To control possible loops even when the Unique Games instance has
parallel constraints with different permutations, first pad its
alphabet from $[\ell]$ to $[\ell]\times[t]$, replacing every
constraint $\pi$ by $(i,j)\mapsto(\pi(i),j)$.
Projecting any satisfying padded label pair to its first coordinate
shows that padding cannot increase the NO optimum. A YES labeling
extends by giving every labeled vertex the same second coordinate.
Regularity and the YES fraction are unchanged.

The AKS edge distribution~\cite[proof of Theorem~3.1]{AKS2011}
first chooses two constraints incident with the same left vertex,
then samples independent coordinate pairs $(x_j,y_j)$ with
probabilities $\Pr[x_j=0,y_j=0]=1-2q$,
$\Pr[x_j=0,y_j=1]=\Pr[x_j=1,y_j=0]=q$, and
$\Pr[x_j=1,y_j=1]=0$. It then permutes the endpoint strings by
the two constraints.

The coordinates $[\ell]\times\{j\}$, for a fixed $j\in[t]$,
form one coordinate block. Put $c=q^2+(1-q)^2<1$. Conditional on any choice of the two
constraints, a loop requires equality of the two permuted strings
in every one of the $t$ coordinate blocks. In each block, full
string equality implies equality at a fixed output coordinate. If
the two permutations select the same input coordinate there, equality
has probability $1-2q$. If they select different input coordinates,
the bits are independent with marginal probability $q$, so equality
has probability $c=q^2+(1-q)^2$. Both bounds are at most $c$.
Indeed, $c-(1-2q)=2q^2>0$.
The coordinate blocks are independent, so the probability of full
string equality, and hence the total loop weight, is at most $c^t$.
Choose a fixed $t$ with $c^t<\varepsilon_{\mathrm{loop}}$, and delete all loops.
This loses at most $\varepsilon_{\mathrm{loop}}$ internal edge weight in any set,
preserves the YES independent set, and does not increase weighted
degrees. Alphabet padding preserves the source promise and
regularity, so the construction above has the stated bounds.
All operations are deterministic, and $t$ is fixed.
\end{proof}

\subsubsection{Vertex unweighting and uniform degree bounds}

Write the padded alphabet size as $R$ and the fixed bias as
$q=a_0/a_1$, with positive integers $a_0<a_1$. For a binary
string $x\in\{0,1\}^R$, let $|x|$ be its
number of coordinates equal to one. The vertex weights in the AKS
construction, indexed by right vertices $v\in V$ and strings $x$,
have the form
\[
 a_{(v,x)}=
 \frac{a_0^{|x|}(a_1-a_0)^{R-|x|}}{|V|a_1^R},
 \qquad v\in V,\quad x\in\{0,1\}^R.
\]
Thus $N_{\mathrm{cl}}=|V|a_1^R$ is polynomial in the Unique Games
instance, since $a_1$ and $R$ are fixed. Replace each original vertex
$u$ by $m_u=N_{\mathrm{cl}}a_u$ independent copies, called
\emph{clones}, labeled
$(u,1),\ldots,(u,m_u)$. These copies form the \emph{clone block} of $u$.
If several source edges join the same pair of vertices, combine their
weights. For every resulting edge $uv$ of weight $w_{uv}$ in $G^-$,
join each clone of $u$ to each clone of $v$ by an edge of weight
$w_{uv}/(m_um_v)$. There are no other edges.
Denote this edge-weighted, vertex-unweighted
graph by $G_1$. It has $N_{\mathrm{cl}}$ vertices, total edge weight
at most one, and weighted degree at most $2/N_{\mathrm{cl}}$ at each vertex.

A clone of $u$ has weighted degree
$\deg_{G^-}(u)/m_u\le2a_u/(N_{\mathrm{cl}}a_u)=2/N_{\mathrm{cl}}$.
The YES independent set becomes one of the same normalized size by
selecting all its clone blocks.

\begin{lemma}[Density after vertex unweighting
{\cite[Section~4, Step~1]{AKS2011}}]
\label{lem:aks-clone-density}
If $G^-$ satisfies~\eqref{eq:aks-loop-free-density}, then its
clone graph $G_1$ satisfies, for every $S'\subseteq V(G_1)$,
\begin{equation}\label{eq:aks-clone-density}
 e_{G_1}(S')\ge\Gamma_\rho(|S'|/N_{\mathrm{cl}})
 -\varepsilon_{\mathrm{dens}}-\varepsilon_{\mathrm{loop}}.
\end{equation}
\end{lemma}

\begin{proof}
For a chosen set $S'\subseteq V(G_1)$, let $f_u$ be the fraction of
clones of $u$ in $S'$. The cited argument forms a random set
$S\subseteq V(G^-)$ by selecting each source vertex $u$ independently
with probability $f_u$. Since $a_u=m_u/N_{\mathrm{cl}}$,
each clone represents source weight $1/N_{\mathrm{cl}}$. The loop
removal in Lemma~\ref{lem:aks-loop-removal} makes the two endpoints
of each edge distinct, so their selection events are independent.
For an edge $uv$ of weight $w_{uv}$, each of the $m_um_v$ clone pairs
has weight $w_{uv}/(m_um_v)$. The set $S'$ contains $m_uf_u$ clones
of $u$ and $m_vf_v$ of $v$, so these clone pairs contribute
$(m_uf_u)(m_vf_v)w_{uv}/(m_um_v)=w_{uv}f_uf_v$. Thus the expected
source weight and expected internal edge weight are
\[
 \mathbb E[a(S)]=\sum_u a_uf_u=\frac{|S'|}{N_{\mathrm{cl}}},
 \qquad
 \mathbb E[e_{G^-}(S)]=\sum_{uv}w_{uv}f_uf_v=e_{G_1}(S').
\]
Applying the source density bound and then convexity of $\Gamma_\rho$
gives
\[
 \begin{aligned}
 e_{G_1}(S')
 &=\mathbb E[e_{G^-}(S)]\\
 &\geq\mathbb E[\Gamma_\rho(a(S))]
       -\varepsilon_{\mathrm{dens}}-\varepsilon_{\mathrm{loop}}\\
 &\geq\Gamma_\rho\!\left(\mathbb E[a(S)]\right)
       -\varepsilon_{\mathrm{dens}}-\varepsilon_{\mathrm{loop}}\\
 &=\Gamma_\rho(|S'|/N_{\mathrm{cl}})
       -\varepsilon_{\mathrm{dens}}-\varepsilon_{\mathrm{loop}}.
 \end{aligned}
\]
This is the claimed bound.
\end{proof}

For the normalized promise, fix a rational $0<\eta<1/10$.
Let $N$ denote the number of output vertices. The positive rational
constant $\zeta_\eta$ below may depend on $\eta$.

\begin{lemma}[Normalized density promise]
\label{lem:normalized-density-source}
There is a choice of $\zeta_\eta$ for which it is UGC-hard to
distinguish simple edge-weighted graphs on $N$
vertices, with positive rational edge weights, total edge weight
at most one, and weighted degree at most $2/N$, satisfying
\begin{description}
\item[YES:] There is an independent set of size at least $(1/2-\eta)N$.
\item[NO:] Every set of at least $\eta N$ vertices has internal edge
 weight at least $\zeta_\eta$.
\end{description}
The reduction is deterministic and polynomial time.
\end{lemma}

\begin{proof}
Set $q=1/2-\eta/4$ and $\rho=-q/(1-q)$. Choose a positive rational
$g$ with $4g<\Gamma_\rho(\eta)$, take
$\varepsilon_{\mathrm{dens}}\leq g$ and
$\varepsilon_{\mathrm{loop}}=g$, and choose $\gamma\leq\eta/4$
small enough for Theorem~\ref{thm:aks-density-source}.
Apply the regular source of Theorem~\ref{thm:regular-ug-source},
then the loop removal and clone construction above.
The YES independent-set density is at least
$q-\gamma\geq1/2-\eta/2\geq1/2-\eta$. For $|S|\geq\eta N$,
equation~\eqref{eq:aks-clone-density} and monotonicity give
\[
 e_{G_1}(S)\geq\Gamma_\rho(\eta)-2g>2g>g.
\]
Thus we may take $\zeta_\eta=g$.

The loop removal makes the clone graph simple after parallel edge
weights are combined. Its total edge weight and degree bounds were
proved above. Rational bias gives rational probabilities in the
finite AKS edge distribution. Enumerating the choices computes the
edge weights exactly. The clone weights are then rational as well.
For fixed $\eta$, all alphabet sizes, denominators of the bias,
and padding factors are constants. The displayed clone count is
therefore polynomial, and all operations are deterministic.
\end{proof}

\subsubsection{Deterministic sparsification of a dense-set promise}

To obtain bounded ordinary degree, we sparsify the weighted graph
and discard edges of very small weight. There will be few enough
edges that discarding them loses little total weight. Every remaining
edge then has a lower weight bound, so a bound on weighted degree
limits the number of incident edges. We use spectral sparsification
to preserve the internal weight of large vertex sets throughout this
process.

For a loopless weighted graph $G$,
let $w_{uv}$ be the total weight of edges between distinct vertices
$u,v$. Its \emph{Laplacian} is the symmetric matrix with
$L_G(u,u)=\deg_G(u)$ and $L_G(u,v)=-w_{uv}$ for $u\ne v$.
Expanding the quadratic form gives
\[
 x^{\mathsf T}L_Gx
 =\sum_{uv\in E(G)}w_{uv}(x_u-x_v)^2.
\]
For $S\subseteq V(G)$, the quadratic form counts boundary edge
weight, so
$\mathbf1_S^{\mathsf T}L_G\mathbf1_S=w(\cut_G(S))$ and
$\mathbf1_{\{v\}}^{\mathsf T}L_G\mathbf1_{\{v\}}=\deg_G(v)$.
The notation $L_G\preceq L_H$ means that
$x^{\mathsf T}L_Gx\le x^{\mathsf T}L_Hx$ for every vector $x$.
A \emph{spectral sparsifier} is a weighted subgraph with fewer
edges whose Laplacian preserves these quadratic forms within
multiplicative bounds. Finally,
$2e_G(S)=\vol_G(S)-w(\cut_G(S))$: an internal edge
contributes twice to the volume, while a crossing edge contributes
once to both the volume and the cut.

\begin{theorem}[BSS spectral sparsification
{\cite[Theorem~1.1]{BSS2012}}]
\label{thm:bss-spectral-sparsification}
For every fixed $q>1$, a loopless undirected weighted graph $G$
on $N\geq1$ vertices has a reweighted spanning subgraph $G'$ with
at most $\lceil q(N-1)\rceil$ edges such that
\[
 L_G\preceq L_{G'}\preceq
 \frac{q+1+2\sqrt q}{q+1-2\sqrt q}L_G.
\]
Such a graph can be constructed deterministically in polynomial time.
\end{theorem}

\begin{lemma}\label{lem:deterministic-density-sparsification}
Fix positive constants $C_0,\zeta,\eta$. Let $G$ be a loopless
edge-weighted graph on $N\ge1$ vertices, with parallel edge weights
already combined, of total edge weight at most one and weighted degree
at most $C_0/N$ at every vertex. Suppose that
every set of at least $\eta N$ vertices has internal edge weight
at least $\zeta$. There is a deterministic polynomial-time
construction of a simple unweighted graph $J$ with
$V(J)=V(G)$ and $E(J)\subseteq E(G)$, whose
maximum degree is bounded by a constant depending only on
$C_0,\zeta$, and whose independent sets all have size less than
$\eta N$.
\end{lemma}

\begin{proof}
Choose a fixed rational $0<\theta\le\min\{1/12,\zeta/14\}$,
and put $q=\theta^{-2}$ and $d=\lceil q\rceil$.
Apply Theorem~\ref{thm:bss-spectral-sparsification} to obtain a
reweighted subgraph $G'$ with at most
$\lceil q(N-1)\rceil\leq dN$ edges and the stated spectral bounds.
Substituting $q=\theta^{-2}$ gives
\[
 \frac{q+1+2\sqrt q}{q+1-2\sqrt q}
 =\left(\frac{1+\theta}{1-\theta}\right)^2
 =1+\frac{4\theta}{(1-\theta)^2}
 \le1+\frac{576}{121}\theta<1+5\theta,
\]
since $\theta\le1/12$.

Applying this sandwich to $\mathbf1_{\{v\}}$ gives
$\deg_G(v)\le\deg_{G'}(v)\le(1+5\theta)C_0/N$. Summing the lower
bounds over $v\in S$ gives
$\operatorname{vol}_{G'}(S)\ge\operatorname{vol}_G(S)$.
Applying the sandwich to $\mathbf1_S$ gives
$w(\cut_{G'}(S))\le(1+5\theta)w(\cut_G(S))$. Every cut has weight at most the total edge weight, so
$w(\cut_G(S))\le1$. Also $\theta\le\zeta/14$ gives
$5\theta/2\leq5\zeta/28\leq\zeta/4$. Therefore, whenever $|S|\ge\eta N$,
\begin{align*}
 2e_{G'}(S)&=\operatorname{vol}_{G'}(S)-w(\cut_{G'}(S))\\
 &\ge\operatorname{vol}_G(S)-(1+5\theta)w(\cut_G(S))\\
 &=2e_G(S)-5\theta w(\cut_G(S)),\\
 e_{G'}(S)&\ge e_G(S)-\tfrac{5\theta}{2}w(\cut_G(S))
 \ge\zeta-\tfrac{5\theta}{2}\ge3\zeta/4.
\end{align*}
Write $w'_e$ for the weight of $e\in E(G')$. Define the unweighted
graph $J$ by $V(J)=V(G)$ and
$E(J)=\{e\in E(G'):w'_e\ge\zeta/(4dN)\}$.
The edges in $E(G')\setminus E(J)$ have total weight at most
$|E(G')|\zeta/(4dN)\le\zeta/4$.
Thus every $S\subseteq V(G)$ with $|S|\ge\eta N$ satisfies
\[
 \sum_{\substack{e\in E(J)\\e\subseteq S}}w'_e
 \ge e_{G'}(S)-\sum_{e\in E(G')\setminus E(J)}w'_e
 \ge\frac{3\zeta}{4}-\frac\zeta4=\frac\zeta2>0.
\]
Consequently $J[S]$ contains an edge.
Thus every independent set of $J$ has size less than $\eta N$.

For $v\in V(J)$, every $e\in E(J)$ incident with $v$ has
$w'_e\ge\zeta/(4dN)$. Hence
\[
 \deg_J(v)\le\frac{4dN}{\zeta}\deg_{G'}(v)
 \le\frac{4(1+5\theta)C_0d}{\zeta},
\]
a constant depending only on $C_0,\zeta$. The construction is
deterministic and polynomial time by
Theorem~\ref{thm:bss-spectral-sparsification}.
\end{proof}

\begin{restatedtheorem}{Theorem~\ref{thm:bounded-degree-kr-source}}
For every fixed rational $0<\eta<1/10$, there is a constant
$d_\eta$ such that the following promise problem is Unique Games-hard.
The input is a simple unweighted graph $G$ on $n\geq1$ vertices
of maximum degree at most $d_\eta$. Distinguish the following cases:
\begin{description}
 \item[YES:] $\alphaI(G)\geq(1/2-\eta)n$.
 \item[NO:] $\alphaI(G)<\eta n$.
\end{description}
The reduction is deterministic and polynomial time.
\end{restatedtheorem}

\begin{proof}[Proof of Theorem~\ref{thm:bounded-degree-kr-source}]
Take the graph from Lemma~\ref{lem:normalized-density-source} and
apply Lemma~\ref{lem:deterministic-density-sparsification} with
$C_0=2$ and $\zeta=\zeta_\eta$. Deleting edges preserves the YES
independent set. In the NO case, every set of at least $\eta N$
vertices still contains an edge, so every independent set has size
less than $\eta N$. The output degree depends only on $\eta$.
Both constructions are deterministic and polynomial time.
\end{proof}


\subsection{The balanced multipartite source}
\label{app:multipartite-normalization}

\begin{restatedtheorem}{Theorem~\ref{thm:ugc-tripartite-cover}}
For every fixed integer $k\geq3$ and fixed rational
$0<\zeta<1/12$, there is a constant $D_{k,\zeta}$ such that the
following promise problem is Unique Games-hard. The input is a
simple unweighted $k$-partite graph $G$ on $N$ vertices, of maximum
degree at most $D_{k,\zeta}$,
with a specified partition into $k$ independent parts each of size
$N/k$.
\begin{description}
 \item[YES:] $\vc(G)\leq(1/2+\zeta)N$.
 \item[NO:] $\vc(G)\geq(1-1/k-\zeta)N$.
\end{description}
The reduction is deterministic and polynomial time.
\end{restatedtheorem}

\begin{proof}[Proof of Theorem~\ref{thm:ugc-tripartite-cover}]
Apply the bounded-degree source of Theorem~\ref{thm:bounded-degree-kr-source} with
$\eta=\zeta/2$. This is an admissible fixed rational parameter,
since $0<\zeta<1/12$ gives $0<\eta<1/24<1/10$.
For a source graph $G$ on $n$ vertices, form $P_k(G)$ as in
Lemma~\ref{lem:multipartite-product}. It has $N=kn$ vertices,
with $n=N/k$ vertices in each of its $k$ independent layers.
The degree identity in the same lemma gives
\[
 \Delta(P_k(G))=(k-1)\Delta(G)
 \leq(k-1)d_{\zeta/2}=:D_{k,\zeta},
\]
which is constant for fixed $k,\zeta$.

In the YES case, $\alphaI(G)\geq(1/2-\eta)n$.
Taking all $k$ copies of a source independent set gives
\[
 \alphaI(P_k(G))\geq k\alphaI(G)
 \geq k(1/2-\eta)n=(1/2-\eta)N.
\]
Taking complements, we obtain
\[
 \vc(P_k(G))=N-\alphaI(P_k(G))
 \leq(1/2+\eta)N\leq(1/2+\zeta)N.
\]
The last inequality uses $\eta=\zeta/2\leq\zeta$.

In the NO case, $\alphaI(G)<\eta n$, and the upper bound from
Lemma~\ref{lem:multipartite-product} yields
\[
 \alphaI(P_k(G))\leq n+(k-1)\alphaI(G)
 <n+(k-1)\eta n.
\]
Subtracting this strict upper bound from $N=kn$ gives
\[
 \vc(P_k(G))
 >kn-n-(k-1)\eta n
 =\left(1-\frac1k-\frac{k-1}{k}\eta\right)N
 \geq\left(1-\frac1k-\zeta\right)N.
\]
Here $(k-1)\eta/k\leq\eta=\zeta/2<\zeta$, so the claimed
non-strict NO threshold follows in particular. The product has no
loops because edges require distinct layer indices, and no parallel
edges because $G$ is simple. Its $k$ is fixed, so its size is
polynomial in that of $G$. Both the independent-set source reduction
and the product construction are deterministic. This proves the
stated balanced, simple, unweighted, bounded-degree promise.
\end{proof}

\subsection{The global SSE partition promise}
\label{app:sse-normalization}

We use the Gaussian notation $\Gamma_\xi$ and $\Psi$ introduced
before Theorem~\ref{thm:aks-density-source}.

We prove the only Gaussian estimate needed here.

\begin{lemma}[Gaussian separation at medium measure]
\label{lem:gaussian-medium-separation}
There are absolute constants $a,\varepsilon_{\mathrm G}>0$ such that, for
$0<\varepsilon<\varepsilon_{\mathrm G}$, $\varepsilon'\in[\varepsilon/2,\varepsilon]$,
and $1/10\le\mu\le1/2$,
$\mu-\Gamma_{1-\varepsilon'/2}(\mu)\ge a\sqrt\varepsilon$.
\end{lemma}

\begin{proof}
Let $X,Z$ be independent standard normal random variables, put
$\xi=1-\varepsilon'/2$, and write
$Y=\xi X+\sqrt{1-\xi^2}\,Z$. Then $X,Y$ are standard normal
with correlation $\xi$. If $\Pr[X\leq t]=\mu\in[1/10,1/2]$,
the threshold $t$ lies in a fixed bounded interval with $t\leq0$.
Moreover, $\mu-\Gamma_\xi(\mu)=\Pr[X\leq t<Y]$.
We may decrease $\varepsilon_{\mathrm G}$ so that
$\varepsilon_{\mathrm G}\leq1$. Then $0\leq\xi<1$ and
$\sigma:=\sqrt{1-\xi^2}\leq1$. Put $t_0=\Psi^{-1}(1/10)$.
The interval $[t-\sigma,t]$ is contained in the fixed interval
$[t_0-1,0]$. Let
$c_*:=\Pr[Z>1]\min_{x\in[t_0-1,0]}\varphi(x)>0$, where $\varphi$
is the standard normal density. On the event
$t-\sigma\leq X\leq t$ and $Z>1$,
we have $X\leq t$ and
\[
 \begin{aligned}
 Y&>\xi(t-\sigma)+\sigma
  =\xi t+(1-\xi)\sigma
  =t+(1-\xi)(\sigma-t)\geq t.
 \end{aligned}
\]
The last inequality uses $t\leq0$ and $\xi\leq1$. The interval
event for $X$ has probability at least
$\sigma\min_{x\in[t_0-1,0]}\varphi(x)$, and it is independent of
$Z>1$. Hence
\[
 \mu-\Gamma_\xi(\mu)
 \geq c_*\sigma
 \geq\frac{c_*}{2}\sqrt{\varepsilon},
\]
where $1-\xi^2=(1-\xi)(1+\xi)\geq1-\xi=\varepsilon'/2
\geq\varepsilon/4$. Taking $a=c_*/2$ proves the lemma.
\end{proof}

We use the weighted graph notation of Section~\ref{sec:sseh} for
the total edge weight $W$, weighted degree $\deg_G$, volume
$\vol_G$, normalized volume $\mu_G$, expansion $\Phi_G$, and cut
weight $w(\cut_G(A))$.

\begin{theorem}[Global SSE partition gap, weighted normalized form {\cite[Theorem~3.5]{RST}}]
\label{thm:sse-source}
There are absolute constants $c>0$ and $\varepsilon_0>0$ such that,
for every fixed integer $q\geq2$ and every fixed
$0<\varepsilon<\varepsilon_0$, the following promise on finite
undirected graphs with positive rational edge weights is SSE-hard.
Parallel edges and loops are allowed, with $W>0$.
\begin{description}
\item[YES:] There is a partition
  $V(G)=V_1\mathbin{\dot\cup}\cdots\mathbin{\dot\cup}V_q$ with
  $\mu_G(V_i)=1/q$ and $\Phi_G(V_i)\leq2\varepsilon$ for every $i$.
\item[NO:] Every $A\subseteq V(G)$ with
  $\mu_G(A)\in[1/10,9/10]$ satisfies
  $w(\cut_G(A))\geq c\sqrt{\varepsilon}\,W$.
\end{description}
\end{theorem}

\begin{proof}[Proof of Theorem~\ref{thm:sse-source}]
Fix $q$ and $\varepsilon$ as in the statement, taking its threshold
$\varepsilon_0\leq\min\{1,\varepsilon_{\mathrm G}\}$. Choose a rational $\rho$ with
$\sqrt{1-\varepsilon/2}<\rho<\sqrt{1-\varepsilon/4}$. Put
$\varepsilon'=2(1-\rho^2)$. Then $\rho$ and $\varepsilon'$ are
rational, and $\varepsilon'\in(\varepsilon/2,\varepsilon)$.
Instantiate Theorem~3.5 of~\cite{RST} with error parameter
$\varepsilon'$ and a fixed rational source error $\gamma$, to be
chosen below.

\emph{YES case.}
In the YES case the source theorem gives $q$ disjoint classes of
normalized volume $1/q$. The proof of RST Theorem~3.5~\cite{RST} uses
$\rho^2=1-\varepsilon'/2$ and, by its Lemma~6.2, bounds every class
expansion by
\[
 2(1-\rho^2+\eta+2\varepsilon_V)
  +(1-\rho^2+\eta+2\varepsilon_V)^2
  +\frac{(1-\rho^2)\beta}{\rho^2}+2^{-\Omega(k)}.
\]
Their parameter choices satisfy
$\varepsilon_V=\varepsilon'^2$, with their small-set source parameter
$\eta\leq\varepsilon'^2$, $\beta\leq\varepsilon'$, and
$k$ is a sufficiently large constant multiple of
$\log(1/\varepsilon')$. Thus their $2^{-\Omega(k)}$ error is
$O(\varepsilon'^2)$, so for an absolute constant $C_{\mathrm{RST}}$,
$2^{-\Omega(k)}\leq C_{\mathrm{RST}}\varepsilon'^2$. Put
$u=1-\rho^2+\eta+2\varepsilon_V$. Since
$1-\rho^2=\varepsilon'/2$, we have
$u\leq\varepsilon'/2+3\varepsilon'^2$. For $\varepsilon'\leq1/6$,
$u\leq\varepsilon'$, so the first two terms in the displayed bound
are at most
\[
 2u+u^2\leq2(\varepsilon'/2+3\varepsilon'^2)+\varepsilon'^2
 =\varepsilon'+7\varepsilon'^2.
\]
Also $\rho^2=1-\varepsilon'/2\geq1/2$, and hence
\[
 \frac{(1-\rho^2)\beta}{\rho^2}
 \leq\frac{(\varepsilon'/2)\varepsilon'}{1/2}
 =\varepsilon'^2.
\]
Thus the full bound is at most
$\varepsilon'+(8+C_{\mathrm{RST}})\varepsilon'^2$. It is at most
$2\varepsilon'$ after reducing the theorem's $\varepsilon_0$ so that
$\varepsilon_0\leq\min\{1/6,(8+C_{\mathrm{RST}})^{-1}\}$. The error constant does
not depend on $q$. The sum of the classes' normalized
volumes is $q\cdot(1/q)=1$. Since normalized volume is additive
on disjoint sets and the whole graph has normalized volume one,
the vertices outside these classes have total volume zero.
Such vertices have degree zero and can be assigned to any class
without changing its volume or its boundary. We therefore obtain
a partition of the entire vertex set with the required
class volumes. Use the rational-parameter version of the finite RST
reduction. Its output graph has rational edge weights. The bound
above gives each class expansion at most
$2\varepsilon'\leq2\varepsilon$.

\emph{NO case.}
In the NO case, Theorem~3.5 gives, for every set of positive volume,
the lower bound
\[
  \Phi_G(A)\geq
  1-\frac{\Gamma_{1-\varepsilon'/2}(\mu_G(A))+\gamma}{\mu_G(A)}.
\]
Let $a$ be the constant in Lemma~\ref{lem:gaussian-medium-separation},
and choose a fixed rational $0<\gamma\le a\sqrt\varepsilon/20$, as the source
theorem permits.
For $\mu_G(A)\in[1/10,1/2]$, multiply the source expansion
inequality by $2\mu_G(A)$. Since
$\Phi_G(A)=w(\cut_G(A))/\vol_G(A)$ and
$\vol_G(A)=2W\mu_G(A)$, this gives
\[
 \begin{aligned}
  \frac{w(\cut_G(A))}{W}
  &=2\mu_G(A)\Phi_G(A)\\
  &\geq2\mu_G(A)-2\Gamma_{1-\varepsilon'/2}(\mu_G(A))-2\gamma\\
  &=2\bigl(\mu_G(A)-\Gamma_{1-\varepsilon'/2}(\mu_G(A))\bigr)-2\gamma\\
  &\geq2a\sqrt{\varepsilon}-2\gamma
   \geq\left(2-\frac1{10}\right)a\sqrt{\varepsilon}
   \geq\frac{a}{10}\sqrt{\varepsilon}.
 \end{aligned}
\]
Here $2\gamma\leq(a/10)\sqrt{\varepsilon}$. If
$\mu_G(A)\in[1/2,9/10]$, its complement has measure
$1-\mu_G(A)\in[1/10,1/2]$, since volume is additive and
$\mu_G(V(G))=1$. A boundary edge has one endpoint in
each side, so $\cut_G(V(G)\setminus A)=\cut_G(A)$ as multisets.
Applying the preceding argument to the complement therefore gives
the same cut bound for $A$. Thus the
claimed NO cut bound holds throughout $[1/10,9/10]$ with $c=a/10$.
All parameters are fixed constants for each source
promise, so the source reduction remains polynomial for each fixed
$(q,\varepsilon)$.
\end{proof}

\subsubsection{Rounding to an unweighted host}

Theorem~\ref{thm:sse-source} gives a weighted multigraph, whereas the
output line graph has one vertex per unweighted host-edge copy. Replace each
weighted edge $e$ by $\lfloor Mw_e/W\rfloor$ parallel copies, where
$W$ is the total source edge weight and $M$ is chosen below.
Rounding introduces an additive error in normalized
volumes and cut sizes, so the source promise uses the wider interval
$[1/10,9/10]$ rather than the interval $[1/4,3/4]$ required by
Lemma~\ref{lem:linegraph-separators}.
The conversion proves the unweighted source promise stated in
Lemma~\ref{lem:sse-unweighted-gap}. The degree reduction is applied
later to its line graph, through incidence clique factors, in
Section~\ref{subsec:bounded-degree-consequences}.

\begin{restatedlemma}{Lemma~\ref{lem:sse-unweighted-gap}}
There are absolute constants $c>0$ and $\varepsilon_0>0$ such that,
for every fixed integer $q\geq2$ and every fixed rational
$0<\varepsilon<\varepsilon_0$, the following promise on unweighted
multigraphs $G$ with $m>0$ edges, allowing loops, is SSE-hard:
\begin{description}
\item[YES:] There is a partition
  $V(G)=V_1\mathbin{\dot\cup}\cdots\mathbin{\dot\cup}V_q$ such that
  $(1-\varepsilon)/q\leq\mu_G(V_i)\leq(1+\varepsilon)/q$ and
  $\Phi_G(V_i)\leq3\varepsilon$ for every $i$.
\item[NO:] Every $A\subseteq V(G)$ with
  $\mu_G(A)\in[1/4,3/4]$ satisfies
  $|\cut_G(A)|\geq(c/2)\sqrt{\varepsilon}\,m$.
\end{description}
\end{restatedlemma}

\begin{proof}[Proof of Lemma~\ref{lem:sse-unweighted-gap}]
Let $c$ and $\varepsilon_0$ be the constants from
Theorem~\ref{thm:sse-source}. Decrease $\varepsilon_0$ so that every
$0<\varepsilon<\varepsilon_0$ satisfies $\varepsilon\leq1/10$ and
$\varepsilon^2\leq(c/2)\sqrt\varepsilon$.
Fix $q\geq2$ and a rational $\varepsilon$ in this range, and set
$\theta=\min\{\varepsilon^2,\varepsilon/(2q)\}$.

\emph{Construction.}
Let $G$ be a weighted source instance with $k\geq1$ edges of weights
$w_e$ and total weight $W$, counting parallel edges separately.
Set $M=\lceil k/\theta\rceil$ and $r_e=\lfloor Mw_e/W\rfloor$.
Form the unweighted multigraph $G^*$ on $V(G)$ by replacing each
edge $e$ by $r_e$ copies with the same endpoints, with no other edges.
A loop is replaced by loops, and an edge with $r_e=0$ is removed.
Write $m=\sum_e r_e$ and output $G^*$.

\emph{Rounding estimates.}
For each input edge $e$, write $\omega_e=Mw_e/W-r_e\in[0,1)$.
Since $k\leq\theta M$ and $\omega_e<1$ for each of the $k$ input edges,
$(1-\theta)M\leq M-k<M-\sum_e\omega_e=m\leq M$.
Thus $m>0$, and $G^*$ has at most $M=O(k/\theta)$ edges.
Since $q,\theta$ are fixed, the reduction has polynomial size.

For a fixed vertex set $A$, an edge contributes zero, one, or two
endpoint incidences to its volume. A loop in $A$ contributes two.
Since $w_e/W-r_e/M=\omega_e/M$, edge $e$ contributes between zero
and $\omega_e/M$ to the difference of normalized volumes. Summing gives
\begin{equation}\label{eq:round-volume}
 0\leq\frac{\vol_G(A)}{2W}
       -\frac{\vol_{G^*}(A)}{2M}
 \leq\frac1M\sum_e\omega_e
 <\frac{k}{M}\leq\theta.
\end{equation}

For the cut, only crossing edges contribute. Thus
\begin{equation}\label{eq:round-cut}
  |\cut_{G^*}(A)|
  =M\frac{w(\cut_G(A))}{W}-\sum_{e\in\cut_G(A)}\omega_e
  \geq M\left(\frac{w(\cut_G(A))}{W}-\theta\right),
\end{equation}
because $\sum_{e\in\cut_G(A)}\omega_e<k\leq\theta M$.

\emph{YES case.}
Use the source partition $(V_1,\ldots,V_q)$.
Since $\mu_G(V_i)=1/q$, equation~\eqref{eq:round-volume} gives
$1/q-\theta\leq\vol_{G^*}(V_i)/(2M)\leq1/q$.
Also, $(1-\theta)M\leq m\leq M$, so
$1-\theta\leq m/M\leq1$. Since
$\mu_{G^*}(V_i)=\bigl(\vol_{G^*}(V_i)/(2M)\bigr)/(m/M)$, it follows
that
\[
\begin{aligned}
 \mu_{G^*}(V_i)
 &\geq1/q-\theta
 \geq(1-\varepsilon/2)/q
 \geq(1-\varepsilon)/q,\\
 \mu_{G^*}(V_i)
 &\leq\frac{1/q}{1-\theta}
 \leq\frac{1+\varepsilon}{q}.
\end{aligned}
\]
Here the last inequality uses
$\theta\leq\varepsilon/(2q)\leq\varepsilon/(1+\varepsilon)$.
The source bounds give
$w(\cut_G(V_i))=\Phi_G(V_i)\vol_G(V_i)
\leq2\varepsilon\cdot2W/q=4\varepsilon W/q$.
Since $r_e\leq Mw_e/W$, summing the rounded edge multiplicities over
the cut gives
$|\cut_{G^*}(V_i)|\leq Mw(\cut_G(V_i))/W\leq4\varepsilon M/q$.
Consequently
\[
  \Phi_{G^*}(V_i)
  \leq\frac{4\varepsilon M/q}{2M(1/q-\theta)}
  =\frac{2\varepsilon}{1-q\theta}
  \leq\frac{2\varepsilon}{1-\varepsilon/2}
  \leq\frac{40}{19}\varepsilon<3\varepsilon.
\]
The penultimate inequality uses $\varepsilon\le1/10$.

\emph{NO case.}
Let $A$ be any set with
$\mu_{G^*}(A)\in[1/4,3/4]$. Since
$\vol_{G^*}(V(G))/(2M)=m/M$, its volume measured at scale $2M$
satisfies $\vol_{G^*}(A)/(2M)=(m/M)\mu_{G^*}(A)
\in[(1-\theta)/4,3/4]$.
The volume bound~\eqref{eq:round-volume} says that this scaled volume
is at most $\mu_G(A)$ and differs from it by at most $\theta$. Hence
\[
 \frac{1-\theta}{4}
 \leq\frac{\vol_{G^*}(A)}{2M}
 \leq\mu_G(A)
 \leq\frac{\vol_{G^*}(A)}{2M}+\theta
 \leq\frac34+\theta.
\]
Since $\theta\le1/10$, the lower endpoint is at least $9/40>1/10$
and the upper endpoint is at most $17/20<9/10$. Thus
$\mu_G(A)\in[1/10,9/10]$. The source NO promise
therefore gives $w(\cut_G(A))/W\geq c\sqrt{\varepsilon}$. Applying
\eqref{eq:round-cut}, $\theta\leq(c/2)\sqrt\varepsilon$, and $M\geq m$ gives
\[
  |\cut_{G^*}(A)|\geq M(c\sqrt\varepsilon-\theta)
  \geq(c/2)\sqrt\varepsilon\,M
  \geq(c/2)\sqrt\varepsilon\,m.
\]
\end{proof}

\subsection{Provenance of the NP-hardness construction}
\label{subsec:np-provenance}

The reduction in Section~\ref{subsec:np-treedepth} passes from a constraint
system to tripartite Vertex Cover and then to Treedepth, following
Bonnet, Neuen, and Soko\l owski~\cite[Sections~3 and~4]{BNS}. Their
truth-block and assignment-witness construction adapts Bonnet's earlier
treewidth reduction~\cite[full version, Section~3.1]{BonnetTreewidth}.
We inherit two opposite truth blocks per variable, complete bipartite
edges between them, and witnesses for satisfying local assignments.
For each variable in its equation, a witness is adjacent to the truth
block assigning the value opposite to the witness's represented value.
Thus a selected truth block is compatible with the witnesses
that agree with its assignment. Variable-dependent block
sizes also occur in Bonnet's full version~\cite[Section~3.5]{BonnetTreewidth}.

Here the source predicates are three-variable parity equations. Each
has four satisfying assignments. We add a star on their four
witnesses, place its monochromatic witness $000$ or $111$ in a truth
part, and place the other three in the third part. We choose the block
size $5d_x/6$ by analyzing arbitrary partial assignments: uniformly
completing the unassigned variables bounds the excess witnesses of an
equation by $5r_j/6$, where $r_j$ is its number of unassigned variables.
This is the soundness argument proved in Lemma~\ref{lem:xor-cover}.
The corresponding argument in~\cite[proof of Lemma~10]{BNS} uses
variable-by-variable exchanges in a vertex cover.
Pairing the source with a complemented copy then makes all three
parts have size $3m$. These are the modifications to the earlier
encoding used in this proof.

The anchor graph is the one in~\cite[Lemma~6]{BNS}. We use their
identity as a black box at its stated anchor length, together
with the cover-first upper bound and Claims~7 and~9 of their proof.
The main-text corollary extends the identity to anchors as short
as the largest part. Its derivation is given in
Appendix~\ref{app:anchor-proof}. For our balanced source graph this
allows $\ell=3m$ even though $\vc(G_F)\ge11m/2$. Combined with the
parity-witness bound, the anchor identity yields the factor below
$18/17$.

\subsection{The parity source normalization}
\label{app:xor-normalization}

\begin{lemma}[Balancing parity right-hand sides]\label{lem:parity-balance}
A system $F_0$ of $m_0$ three-variable parity equations can be paired
with a system $F_0'$ on fresh variables so that their union has
$2m_0$ equations, optimum $2\OPT(F_0)$, and $m_0$ equations with
each right-hand side.
\end{lemma}

\begin{proof}
Make a copy $F_0'$ on fresh
variables, replacing every equation $x+y+z=a$ by
$x'+y'+z'=1-a$. For an assignment $\sigma$ to $F_0$, assign
$x'=1-\sigma(x)$ in $F_0'$. Since each equation contains three variables,
\[
 (1-\sigma(x))+(1-\sigma(y))+(1-\sigma(z))
 \equiv1-(\sigma(x)+\sigma(y)+\sigma(z))\pmod2.
\]
Thus the copied equation is satisfied if and only if the original
one is satisfied. Complementing all variable values is a bijection
between assignments of the two systems, so
$\OPT(F_0')=\OPT(F_0)$.

The systems have disjoint variable sets. Any assignment to their
union satisfies at most $\OPT(F_0)+\OPT(F_0')$ equations. Choosing
an optimal assignment independently in each system attains this
bound. Hence their union has $2m_0$ equations and optimum
$2\OPT(F_0)$. Every original equation is paired with one having
the opposite right-hand side, so each right-hand side occurs in
$m_0$ equations.
\end{proof}

\begin{restatedtheorem}{Theorem~\ref{thm:xor-source}}
For every fixed $0<\epsilon<1/4$, there is a constant
$b_\epsilon\geq6$ depending only on $\epsilon$ such that the following
promise problem is NP-hard. The input is an unweighted system $F$ of
$m\geq1$ equations $x+y+z=a\pmod2$, where $a\in\{0,1\}$ and
$x,y,z$ are distinct Boolean variables. Equations may be repeated.
Every variable occurs at most $b_\epsilon$ times, with at most two
possible occurrence counts, and the instances satisfy
\[
 d_x\in6\mathbb Z_{>0},\qquad
 \sum_xd_x=3m,\qquad
 |\{j:a_j=0\}|=|\{j:a_j=1\}|=m/2.
\]
Here $d_x$ is the number of equations containing $x$, and $\OPT(F)$
is the maximum number of equations satisfied by one assignment,
counting repetitions.
\begin{description}
 \item[YES:] $\OPT(F)\geq(1-\epsilon)m$.
 \item[NO:] $\OPT(F)\leq(1/2+\epsilon)m$.
\end{description}
\end{restatedtheorem}

\begin{proof}[Proof of the normalized form in Theorem~\ref{thm:xor-source}]
Hazan, Safra, and Schwartz state the required source gap as
Theorem~10 of~\cite{HazanSafraSchwartz2003}, attributing it to
H\aa stad. Their theorem gives unweighted Max-3-LIN-2 instances in
which each variable occurs a constant number of times, chosen from
two possible values, and no variable repeats within an equation.
Thus the three variables in each source equation are distinct. Apply
their theorem with parameter $\epsilon/2$. Its YES
threshold $1-\epsilon/2$ is at least $1-\epsilon$, and its NO bound
$<1/2+\epsilon/2$ implies the required bound $\leq1/2+\epsilon$.
We use this form directly. Only the balance of the right-hand sides
and the divisibility of occurrence counts need normalization.

We impose the additional conditions in~\eqref{eq:xor-normalization}.
Call the initial system $F_0$, and let $m_0$ be its number of
equations. First delete unused variables. This does not change the
number of equations or the optimum. Apply Lemma~\ref{lem:parity-balance} to pair $F_0$ with its
complemented copy.

Next repeat every equation in the union six times, and call the
resulting system $F$. An assignment either satisfies all six copies
of an equation or satisfies none. Consequently
\[
 m=12m_0,\qquad \OPT(F)=12\OPT(F_0),\qquad
 \frac{\OPT(F)}m=\frac{\OPT(F_0)}{m_0}.
\]
In particular, a YES instance of the original promise satisfies
$\OPT(F)\geq(1-\epsilon)m$, and a NO instance satisfies
$\OPT(F)\leq(1/2+\epsilon)m$.

Each remaining variable occurred at least once before repetition,
and its final occurrence count is six times a positive integer.
The repetition also multiplies the occurrence upper bound by six,
so the normalized bound $b_\epsilon$ still depends only on $\epsilon$.
Thus $d_x\in6\mathbb Z_{>0}$. Counting the three variable
occurrences in every equation gives $\sum_xd_x=3m$.
The repetition creates $6m_0=m/2$ equations with each right-hand
side, so the required balance also holds and $m$ is even. These
operations increase the number of equations by the fixed factor
twelve and therefore take polynomial time. This proves all the
conditions in~\eqref{eq:xor-normalization}.
\end{proof}

\subsection{A cubic independent-set source from the parity gap}
\label{app:cubic-independent-set-source}

We derive the threshold promise in
Theorem~\ref{thm:source-cubic-independent-set} from the existing
parity source. The degree normalization below has an explicitly
computed additive effect on the independence number, so the
threshold remains part of the polynomial-time reduction.

\begin{restatedtheorem}{Theorem~\ref{thm:source-cubic-independent-set}}
There is a fixed rational constant $\delta_0\in(0,1)$ such that the
following promise problem is NP-hard. The input is a simple
unweighted cubic graph $J$ on $n\geq4$ vertices and an integer
$1\leq K\leq n$.
\begin{description}
 \item[YES:] $\alphaI(J)\geq K$.
 \item[NO:] $\alphaI(J)\leq(1-\delta_0)K$.
\end{description}
\end{restatedtheorem}

\begin{proof}[Proof of Theorem~\ref{thm:source-cubic-independent-set}]
Use Theorem~\ref{thm:xor-source} with the fixed error
$\epsilon=1/12$, and construct its witness graph $G_F$ from
Section~\ref{subsec:np-treedepth}. Put $n_0=|V(G_F)|=9m$ and
$e_0=|E(G_F)|$, and take a fixed integer
$D\geq3b_{1/12}$. The construction has maximum degree at most $D$,
and Lemma~\ref{lem:xor-cover} gives
\[
 \alphaI(G_F)=\frac{5m}{2}+\OPT(F).
\]
Consequently its YES independence number is at least $41m/12$,
and its NO independence number is at most $37m/12$. The graph has
no isolated vertices. Every truth vertex meets its opposite truth
block, and every witness meets a truth block.

For each vertex $v$ of degree $d_v\geq1$, replace $v$ by a path
on $2d_v-1$ new vertices. Designate its $d_v$ odd-position vertices
as ports, assign the edges incident with $v$ to distinct ports, and
replace each original edge by an edge between its assigned ports.
Call the resulting simple graph $Q$. Each port has at most two
path neighbors and one external neighbor, so $Q$ has maximum degree
at most three.

An independent set in this path has at most $d_v$ vertices, and it
has $d_v$ only when it contains all the ports. Thus the original
vertices whose paths contribute $d_v$ form an independent set in
$G_F$. Every other path contributes at most $d_v-1$. Conversely,
an independent set in $G_F$ gives one in $Q$ by selecting all ports
in its paths and all even-position vertices in the other paths.
The external edges do not join two selected vertices. These two
bounds prove
\[
 \alphaI(Q)=\alphaI(G_F)+\sum_v(d_v-1)
           =\alphaI(G_F)+2e_0-n_0.
\]
The path construction also gives
\[
 |V(Q)|=4e_0-n_0,\qquad |E(Q)|=5e_0-2n_0.
\]
Hence the number of missing degrees needed to make $Q$ cubic is
\[
 r=\sum_{u\in V(Q)}(3-\deg_Q(u))
   =3|V(Q)|-2|E(Q)|=2e_0+n_0.
\]

For each missing degree at a vertex $u$ of $Q$, take a fresh copy
of $K_4$, subdivide one edge $ab$ with a new vertex $z$, and add
the edge $uz$. All five vertices in this copy are new. The four
original vertices have degree three, and $z$ has degree three
after attachment. The resulting graph $J$ is therefore simple,
unweighted, and cubic.

Each five-vertex attachment has independence number two. Without
$z$, it is $K_4$ with the edge $ab$ removed, whose maximum
independent set is $\{a,b\}$. With $z$ selected, neither $a$ nor
$b$ can be selected, and the two other vertices are adjacent.
Moreover, the pair $\{a,b\}$ avoids the attachment vertex $z$, so
it can be selected whether or not $u$ is selected. Thus the
attachments add exactly $2r$ to the independence number, giving
\[
 \begin{aligned}
  |V(J)|&=|V(Q)|+5r=14e_0+4n_0,\\
  \alphaI(J)&=\alphaI(G_F)+s,
  \qquad s=2e_0-n_0+2r=6e_0+n_0.
 \end{aligned}
\]

Define the integer threshold and the fixed rational gap by
\[
 K=s+\left\lceil\frac{41m}{12}\right\rceil,
 \qquad \delta_0=\frac{1}{81D+41}.
\]
In the YES case, integrality of the independence number gives
$\alphaI(J)\geq K$. In the NO case,
\[
 \alphaI(J)\leq s+\frac{37m}{12}
              \leq K-\frac m3.
\]
Since $e_0\leq Dn_0/2$, $n_0=9m$, and $m\geq1$,
\[
 K\leq6e_0+9m+\frac{41m}{12}+1
   \leq\left(27D+\frac{161}{12}\right)m.
\]
It follows that
\[
 \delta_0K\leq
 \frac{27D+161/12}{81D+41}\,m<\frac m3,
\]
and therefore the NO case satisfies
$\alphaI(J)\leq(1-\delta_0)K$.
The threshold is positive and at most the output order, since
\[
 |V(J)|-K\geq8e_0+3n_0-\frac{41m}{12}-1
            =8e_0+\frac{283m}{12}-1>0.
\]
The bound on $D$ is fixed, so $e_0=O(m)$ and every construction
above has $O(m)$ vertices and edges. All operations and the
threshold computation take deterministic polynomial time.
The source hardness in Theorem~\ref{thm:xor-source} therefore
proves the claimed cubic promise.
\end{proof}
